\begin{proposition}[G-14: First-jet dominance and a stronger padding boundary]\label{G-14}
Let $m\geq1$, and write a quaternary form as
\[
 q(x_0,x_1,x_2,x_3)=\sum_{j=0}^m x_0^j q_j(x_1,x_2,x_3).
\]
The first-jet map
\[
 j_1:\mathcal D_{m,4}\longrightarrow
   \operatorname{Sym}^m\mathbb C^3\oplus
   \operatorname{Sym}^{m-1}\mathbb C^3,
 \qquad q\longmapsto(q_0,q_1),
\]
has Zariski-dense image.

Consequently let $n>m\geq1$, put $a=n-m$, and let $f$ be
any degree-$m$ form in the common ambient GCT setting.
For every $d\geq1$ and every partition $\lambda\vdash nd$
with at most four rows and
\begin{equation}\label{G14:boundary}
 \lambda_1\leq ad+1,
\end{equation}
one has
\[
 m_\lambda(\ell^a f)\leq m_\lambda(\det_n).
\]
If $\lambda_1<ad$, the multiplicity on the left is zero.
Thus every positive multiplicity gap in this sector must
satisfy $\lambda_1\geq ad+2$.
\end{proposition}

\paragraph{Scope and background.}
The proof uses the coefficient grading and first-four-variable
highest-weight restriction of G-05, together with elementary
cyclic cofactors and the finite geometric-product identity
already proved in R-16. It requires characteristic zero.
The jet-dominance calculation is given in full; no higher-jet
dominance or Lefschetz assertion is assumed.

\begin{proof}
The case $m=1$ of the first assertion is immediate. Suppose
$m\geq2$, choose a primitive $m$th root of unity $\zeta$,
and let $A$ be the ternary cyclic pencil with diagonal
entries
\[
 \ell_i=x_1-\zeta^i x_2,\qquad 0\leq i<m,
\]
and entries $x_3$ in positions $(i,i+1)$, with indices
modulo $m$. Its determinant has the two perfect matchings
and equals
\[
 x_1^m-x_2^m+(-1)^{m-1}x_3^m.
\]
We claim that the cofactors of $A$ span every ternary form
of degree $m-1$.

For each $0\leq k\leq m-1$, the cofactors using exactly
$k$ cycle entries are, up to nonzero signs, $x_3^k$ times
the $m$ cyclic translates of a consecutive product of
$h=m-1-k$ of the binary forms $\ell_i$. This follows from
the unique matching after deleting the relevant row and
column. A consecutive product, after a cyclic shift,
has the form
\[
 F_h(x_1,x_2)=\prod_{i=0}^{h-1}(x_1-\zeta^i x_2).
\]
Every coefficient in it is nonzero. Indeed its coefficient
of $x_1^{h-r}x_2^r$ is
\[
 (-1)^r\zeta^{r(r-1)/2}
   \prod_{j=1}^r\frac{1-\zeta^{h-r+j}}{1-\zeta^j},
 \qquad 0\leq r\leq h,
\]
with the empty product interpreted as one. Every exponent
in a nonempty numerator or denominator lies between $1$
and $h<m$. Thus no factor is zero. The cyclic substitution
$x_2\mapsto\zeta x_2$ has distinct characters on the
monomials with $0\leq r\leq h$. Fourier projection of the
$m$ translates therefore spans every binary form of degree
$h$. Varying $k$ proves the claim, since
\[
 \operatorname{Sym}^{m-1}\mathbb C^3
 =\bigoplus_{k=0}^{m-1}x_3^k
        \operatorname{Sym}^{m-1-k}\langle x_1,x_2\rangle.
\]

Now consider the parameter map
\[
 (A',B)\longmapsto j_1\det(A'+x_0B),
\]
where $A'$ is an arbitrary ternary linear pencil and $B$
is a constant $m\times m$ matrix. At the point $(A,0)$,
the derivative in the $A'$ directions changes $q_0$ by
arbitrary linear-form multiples of the cofactors of $A$.
By the claim these span all of
$\operatorname{Sym}^m\mathbb C^3$. It does not change
$q_1$, since the coefficient of $x_0$ is linear in $B$.
The derivative in the $B$ directions changes $q_1$ by
arbitrary constant linear combinations of the cofactors,
which span $\operatorname{Sym}^{m-1}\mathbb C^3$, and
does not change $q_0$. The differential is consequently
surjective onto the direct sum. The characteristic-zero
Jacobian criterion for dominance proves the first assertion.

For the multiplicity consequence, work first in four
variables and use the positive coefficient weight convention
of G-05/G-11 (equivalently, dualize the coordinate-ring
representations consistently). In a polynomial irreducible
of shape $\lambda$, every weight has each coordinate at
most $\lambda_1$. For example, the Schur construction uses
$\lambda_1$ exterior-power columns, and a given basis vector
can occur at most once in each column. This proves the
required bound on the $x_0$ coordinate of every weight.

Let $R=\mathbb C[\operatorname{Sym}^n\mathbb C^4]$ be
the coefficient polynomial ring, and let $R_{d,\lambda}$
be its degree-$d$ $\lambda$-isotypic subspace. Take
\[
 P\in I(\mathcal D_{n,4})\cap R_{d,\lambda}.
\]
The weight bound says that every coefficient monomial
occurring in $P$ has total $x_0$ weight at most
$\lambda_1$. Evaluate $P$ at a form $x_0^a q$. A
coefficient coming from $q_j$ has $x_0$ weight $a+j$;
a product of $d$ coefficients therefore has weight
\[
 ad+j_1+\cdots+j_d.
\]
All coefficient variables of smaller individual $x_0$
weight vanish on $x_0^a q$. If $\lambda_1<ad$, nothing
survives and every polynomial in $R_{d,\lambda}$ vanishes
on every such padded form.

Under \eqref{G14:boundary}, the restriction
$P(x_0^a q)$ depends only on $q_0,q_1$, because every
surviving coefficient monomial has
$j_1+\cdots+j_d\leq1$. It is a polynomial function of
these first-jet data. For every $q\in\mathcal D_{m,4}$,
block padding gives $x_0^a q\in\mathcal D_{n,4}$,
including the closure points. Hence this jet polynomial
vanishes on the dense image of $j_1$ just proved, and
therefore vanishes for every pair $(q_0,q_1)$. It follows
that $P(x_0^a q)=0$ for every quaternary form $q$ of
degree $m$.

To allow any padding form $L$, use a $\mathrm{GL}_4$
change taking $L$ to $x_0$. The determinant ideal and
the isotypic subspace are both stable under that group,
so the preceding conclusion applies to the corresponding
translate of $P$. Thus $P$ vanishes on every $L^a q$,
and in particular on the restricted padded orbit closure
associated with $\ell^a f$.
We have proved the containment of the determinant-ideal
$\lambda$-isotypic part in the corresponding padded ideal
part. Passing to the characteristic-zero semisimple
degree-$d$ quotients gives the stated multiplicity
inequality. The first-four-variable highest-weight
restriction identifies these multiplicities with those
in the full ambient orbit closures, as in G-05/R-14.
The same argument gives the zero assertion when
$\lambda_1<ad$.
\end{proof}

\paragraph{A precisely stated higher-jet continuation.}
For $0\leq h\leq m$, suppose that the projection
\[
 j_h:\mathcal D_{m,4}\longrightarrow
   \bigoplus_{j=0}^h\operatorname{Sym}^{m-j}\mathbb C^3
\]
has dense image. The identical coefficient-weight argument
then proves the same multiplicity inequality whenever
$\lambda_1\leq ad+h$: every surviving coefficient
monomial depends only on those jets. This is a conditional
continuation with an explicit geometric hypothesis.
The first-jet case above is proved without that hypothesis.
No dominance assertion for $h=2$ or $h=3$ is made here.
